\begin{figure}
    \centering
    \begin{tikzpicture}[->,>=stealth',shorten >=1pt,right,node distance=1.6cm,semithick,
        index/.style={
        rectangle,
        draw=black,
        thick,
        fill=blue!10,
        text width=4cm,
        align=center,
        rounded corners,
        minimum height=2em
        },
        query/.style={
        rectangle,
        draw=black,
        thick,
        fill=ForestGreen!15,
        text width=4cm,
        align=center,
        rounded corners,
        minimum height=2em
    }]
      \node[index] (Source Text Documents) {Source Documents};

      \node[index] (Source Text Chunks) [below of=Source Text Documents] {Text Chunks};
      \path (Source Text Documents) edge node[text width=2cm] {\small text extraction and chunking} (Source Text Chunks);
      
      \node[index] (Graph Element Instances) [below of=Source Text Chunks] {Entities \& Relationships};
      \path (Source Text Chunks) edge node[text width=2.5cm] {\small domain-tailored summarization} (Graph Element Instances);
      
      \node[index] (Graph Element Summaries) [below of=Graph Element Instances] {Knowledge Graph};
      \path (Graph Element Instances) edge node[text width=2.5cm] {\small domain-tailored summarization} (Graph Element Summaries);
      
      \node[index] (Graph Community Structures) [right=3cm of Graph Element Summaries] {Graph Communities};
      \path (Graph Element Summaries) edge node[text width=2cm,auto,align=center] {\small community detection} (Graph Community Structures);
      
      \node[index]  (Graph Community Summaries) [above of=Graph Community Structures] {Community Summaries};
      \path (Graph Community Structures) edge node[text width=2.5cm] {\small domain-tailored summarization} (Graph Community Summaries);
      
      \node[query]  (Community Query Summaries) [above of=Graph Community Summaries] {Community Answers};
      \path (Graph Community Summaries) edge node[text width=2.5cm] {\small query-focused summarization} (Community Query Summaries);
      
      \node[query] (Global Query Summary) [above of=Community Query Summaries] {Global Answer};
      \path (Community Query Summaries) edge node[text width=2.5cm] {\small query-focused summarization} (Global Query Summary);

      \node[index, font=\itshape] (Indexing Time) [below of=Graph Element Summaries] {Indexing Time};

      \node[query, font=\itshape] (Query Time) [below of=Graph Community Structures] {Query Time};

      \node (Stage) [right=0.3cm of Indexing Time, font=\bfseries] {Pipeline Stage};
      
    \end{tikzpicture}
    \caption{
    Graph RAG pipeline using an LLM-derived graph index of source document text.
    This graph index spans nodes (e.g., entities), edges (e.g., relationships), and covariates (e.g., claims) that have been detected, extracted, and summarized by LLM prompts tailored to the domain of the dataset.
    Community detection (e.g., Leiden,~\citeauthor{Traag2019Leiden}, \citeyear{Traag2019Leiden}) is used to partition the graph index into groups of elements (nodes, edges, covariates) that the LLM can summarize in parallel at both indexing time and query time.
    The ``global answer'' to a given query is produced using a final round of query-focused summarization over all community summaries reporting relevance to that query.
    }
    \label{fig:pipeline}
\end{figure}
